\section{Absorbed by the discourse}
\label{sec:15.6}
The third instance is the discourse that examines the machine, and this argument is inside it. Objections raised there became job titles, benchmarks and certifications, and nobody outside the industry holds them.

The argument that alignment and AI ethics came back as something the industry could use is older than the current wave and was made from inside: Ochigame's 2019 piece arguing that “ethical AI” discourse was constructed to forestall legally enforceable restriction, that industry investment in academic ethics was investment in a \emph{substitute} for regulation, and that the substitute works by converting political questions into technical ones a company can then be trusted to answer \autocite{ochigame2019invention}.

\runin{Where the objections went}

An objection raised inside a lab becomes a paper or a resignation, and the labs collecting them can point to the collection as evidence of seriousness. One stage on, alignment is the job title and the voluntary commitment is the certification. Both were somebody's objections first, and each is now a criterion an objector is measured against or a role somebody else holds. There's no route by which a researcher's conclusion that a system shouldn't ship becomes a decision that it doesn't, because the lab that decides is the one that pays.

All four operations are present in the discourse, and none needs a movement taking power. What a movement could do with them is the danger.

\runin{The remedy}

The structural remedy is a refusal whose cost the refuser can see and choose to bear, and a body with standing to receive it — an external authority that can stop a deployment, to which an objection can be filed by someone who doesn't depend on the objection being welcome. A decade of the voluntary version hasn't produced it. Until something of that kind exists, everything published under the heading of alignment is collected by a channel whose owners decide what it moves.

\runin{A transfer hollowed in public}

Voting-rights litigation after 2026 is a worked case of measuring the trend, not the level (\S\ref{sec:14.4}). \emph{Louisiana v. Callais} left Section~2 of the Voting Rights Act in force and raised what a challenger must prove: when a map is defended as partisan, the challenger has to disentangle race from politics and show a strong inference that the state drew its lines because of race \autocite{callais2026}. Objectors kept standing to invoke the rule; what moved was the burden of using it, and it moved onto the party that doesn't hold the evidence, since the state draws the map and keeps the record of how it was drawn. That is neither recuperation nor open repeal. It is a transfer hollowed in public, by a reasoned decision that presents itself as preserving the rule it narrows. The level is a statute unchanged since 1982; the trend is the cost of bringing a claim under it.
